\subsection{Compilation into finite numerical tables}\label{sec:effective}
A covering theorem supplies a finite set, but does not in general supply
an algorithm for choosing that set. Moreover, a table with finitely many
entries may still contain noncomputable real numbers. We separate these
issues. The construction below uses finite grids of feasible histories,
ordinary rational comparisons, and approximations to a finite moment
table. Neither rank detection nor a decision whether two exponents
coincide is among its operations.

\subsubsection{A compilation theorem for compact state systems}
Let $T\ge1$ be fixed, let $U=[\eta,1-\eta]^J$, with rational
$0<\eta<1/2$, and let $\mathcal X$ be finite. For $0\le n\le T$
let $S_n\subset[-B,B]^{d_n}$ be compact, with $S_0$ a singleton.
There are transition maps
\[
 T_{n,x}:S_n\times U\longrightarrow S_{n+1},\qquad
 S_{n+1}=\bigcup_{x\in\mathcal X}T_{n,x}(S_n,U).
\]
Norms in this section are maximum norms. Suppose
\begin{equation}\label{eq:compiler-lipschitz}
 \|T_{n,x}(s,u)-T_{n,x}(s',u')\|_\infty
 \le L_n\|s-s'\|_\infty+G_n\|u-u'\|_\infty.
\end{equation}
Each stage has finitely many queries, with a prediction map
$Q_n:S_n\to[0,1]^{b_n}$ of Lipschitz constant $K_n$.
The bounds and all dimensions are fixed, independent of the state budget.
Write $H_n(w,u_1,\ldots,u_n)$ for the exact state of the report word
$w\in\mathcal X^n$. Iteration of \eqref{eq:compiler-lipschitz} shows that
its command dependence is $A_n$-Lipschitz, where
\[
 A_0=0,\qquad A_{n+1}=L_nA_n+G_n.
\]
For $M\ge1$, define the unrestricted-centre covering radius
\[
 e_n(M)=\inf_{z_1,\ldots,z_M\in\mathbb R^{d_n}}
                  \sup_{s\in S_n}\min_i\|s-z_i\|_\infty.
\]
No spectral representation or dimension formula for $e_n(M)$ is assumed.

\begin{definition}[Finite evaluation data]\label{def:finite-data}
Given a rational grid history and a positive rational tolerance $\tau$,
an evaluation procedure returns a dyadic vector within $\tau$ of its
exact state. It also returns dyadic query predictions within $\tau$ of
the exact predictions at any chosen representative history. This is an
offline procedure. In its effective version it is a specified algorithm
with access to the numerical input oracles of the experiment. In a
noneffective version the finitely many approximations are supplied data.
Neither version gives an oracle to the running transducer.
\end{definition}

Let $U_h\subset U$ be a finite rational $h$-net in the maximum norm,
with $|U_h|\le C_U(1+h^{-1})^J$. At runtime an input includes an index
of any $\bar u\in U_h$ satisfying $\|u-\bar u\|_\infty\le h$.
The guarantee below holds for every such choice. In particular it does
not require deciding equality of an arbitrary real input to a rounding
boundary. An input approximation, followed by rounding that rational
approximation and clipping to $U$, supplies this interface with a fixed
constant adjustment of $h$. For comparison with a memory lower bound,
input coding is a fixed memoryless Borel function of the current command,
or uses fresh independent rounding coins. It is not an additional
history-dependent channel. The upper estimate itself holds for every
name meeting the error bound.

\begin{lemma}[Greedy covering from approximate samples]\label{lem:finite-greedy}
Let $S$ be a compact subset of a normed space. Suppose a finite indexed
list of points $s_j\in S$ is an $a$-net of $S$, and
$\|\widetilde s_j-s_j\|\le\tau$. Farthest-first selection of at most
$M$ points from the rational list $(\widetilde s_j)$, with fixed tie
breaking, selects indices $I$ for which
\[
 \sup_{s\in S}\min_{i\in I}\|s-s_i\|
                 \le 2e_S(M)+a+4\tau.
\]
If there are at most $M$ indices, select all of them. Distinct indices
may have identical approximate vectors.
\end{lemma}
\begin{proof}
Let $r$ be the final covering radius of the approximate list. When the
list has more than $M$ indices, the selected $M$ points and a farthest
point have mutual distances at least $r$. Every cover by $M$ balls of
radius $b$ must put two of these $M+1$ points in the same ball, and
therefore $r\le2b$. This also handles $r=0$. A cover of $S$ with
radius $e_S(M)+\epsilon$ covers the approximate list with radius
$e_S(M)+\epsilon+\tau$, allowing its centres to lie outside $S$.
Thus $r\le2e_S(M)+2\tau$ after $\epsilon\downarrow0$.
For a given $s\in S$, choose $s_j$ within $a$, and then a selected
approximate point within $r$ of $\widetilde s_j$. The triangle
inequality gives $a+2\tau+r$. The case of a list of size at most $M$
is immediate. This is the classical farthest-first argument
\cite{Gonzalez}; the approximation allowance is recorded explicitly.
\end{proof}

\begin{theorem}[Effective compilation without loss of covering order]
\label{thm:compact-compiler}
Under \eqref{eq:compiler-lipschitz} and Definition~\ref{def:finite-data},
there is a finite procedure producing stage-dependent integer transition
tables and dyadic output tables, with at most $M$ state indices at each
stage, such that, on every input-report history,
\begin{equation}\label{eq:compiler-bound}
 \|\widehat Q_n-Q_n(s_n)\|_\infty
 \le C_T\left(\max_{1\le j\le n}e_j(M)+h+\tau\right).
\end{equation}
The running procedure reads only its previous index, the current finite
input code and report, and, at a checkpoint, the new query label. It
stores no numerical state vector, representative history, or oracle.
The constant depends only on the displayed bounds and the fixed horizon.
The synthesis procedure does not need $e_n(M)$ or a covering certificate
as input.
\end{theorem}
\begin{proof}
Enumerate all pairs of a report word in $\mathcal X^n$ and a command
word in $U_h^n$. They are feasible histories by the hypotheses on the
state system. Their exact images form an $A_nh$-net of $S_n$.
Evaluate every image with error at most $\tau$ and apply
Lemma~\ref{lem:finite-greedy}. Associate each selected index with the
\emph{exact} state $c_{n,i}$ of its selected history, for analysis only.
The stored approximations $\widetilde c_{n,i}$ have error at most
$\tau$. The exact representatives cover $S_n$ with radius
\begin{equation}\label{eq:compiled-cover}
 R_n\le2e_n(M)+A_nh+4\tau.
\end{equation}

For each representative history at stage $n$, each $\bar u\in U_h$
and each $x$, append $(\bar u,x)$ to that history and evaluate the
resulting exact state
$y=T_{n,x}(c_{n,i},\bar u)$ with error at most $\tau$.
Choose the next index by minimizing the maximum-norm distance between
this approximation and the stored approximations at stage $n+1$.
All comparisons are between rational numbers. If $j$ is selected, the
triangle inequality, applied to a nearest exact centre, gives
\[
 \|y-c_{n+1,j}\|_\infty\le R_{n+1}+4\tau.
\]
This specifies one integer for every entry of a finite transition table.
The offline representative history is not a runtime variable. In
particular no update is applied to an approximate vector outside the
reachable set: $y$ is defined and evaluated through an actual history.

Let $E_n=\|s_n-c_{n,i_n}\|_\infty$, where $s_n$ denotes the actual
uncompressed state, again only in the analysis. For any admitted finite
input code,
\begin{equation}\label{eq:digital-recurrence}
 E_0=0,\qquad
 E_{n+1}\le L_nE_n+G_nh+2e_{n+1}(M)+A_{n+1}h+8\tau.
\end{equation}
Induction over the fixed horizon bounds $E_n$ by a constant times the
right-hand quantity in \eqref{eq:compiler-bound}. Store a dyadic
approximation to $Q_n(c_{n,i})$ with error at most $\tau$ and clip it
to $[0,1]$. Clipping cannot increase its error to a probability.
The query error is at most $K_nE_n+\tau$, proving the theorem.
There are finitely many evaluations and comparisons, so the procedure
terminates whenever the evaluation procedure does. Nonuniqueness of
representatives or ties does not require an equality test on exact data.
\end{proof}

\begin{proposition}[Size of the finite program]\label{prop:compiler-size}
Put $V=|U_h|$ and $X=|\mathcal X|$. The transition table has at most
$TMVX$ entries, each taking $\lceil\log_2(M+1)\rceil$ bits. If outputs
are written to $b$ binary places, the output tables take at most
$M\sum_nb_n(b+1)$ bits, in addition to fixed metadata. Candidate
enumeration uses at most $\sum_{n=0}^T(XV)^n$ histories. With cached
candidate approximations, greedy selection uses
$O(M\sum_{n=0}^T(XV)^n)$ fixed-dimensional distance operations, and
transition selection uses $O(TM^2VX)$ such operations.
\end{proposition}
\begin{proof}
The transition arguments are the stage, current representative, finite
command code and report. Counting their combinations proves the first
assertion. A query table has one $b$-place probability per representative
and query. At stage $n$ there are $(XV)^n$ grid histories. Maintain for
each candidate its smallest distance to a chosen centre; one new centre
updates this array with one pass, repeated at most $M$ times. For each
transition, comparison against at most $M$ next centres gives the last
count. Exact integer comparison, maximum and subtraction suffice.
These counts exclude the running time of an arbitrary numerical oracle.
\end{proof}

With a supplied clock the only persistent field has at most $M$ values.
An autonomous machine can retain the pair $(n,i_n)$, using at most
$1+(T+1)M$ values. The current input code, table address and returned
output are transient data. They require $O(\log M+\log V+b+\log(T+1))$
bits for fixed dimensions. The finite read-only program is a separate,
explicitly counted resource; it is not claimed to fit into $\log M$ bits.

The fixed fine-mesh realization below is retained in full. Its sufficient
precision and program bounds are strengthened, without changing its
construction or converse, in Section~\ref{sec:certified-resources}.

\subsubsection{Finite data for positive monomial experiments}
Return to the experiment of Section~\ref{sec:collision-geometry}, with
$I=[0,1]$, and put $T=N-1$. Use the full formal raw moment vector
$v_n=(\E[t^{\alpha\cdot a}\mid h])_{|\alpha|=N-n}$, including its
constant entry. Coincident labels remain present. Its dimension is
bounded by a constant depending only on $r,N$.

Choose a fixed bound
\[
 B_f\ge1+\sup_{g,x}\sum_{i=0}^{r-1}|f_i(g,x)|,
 \qquad f_{g,x}(t)=\sum_i f_i(g,x)t^{a_i}.
\]
This uses the coefficient realization, not exponent separation.
For two reachable raw vectors, the update numerator and denominator
have differences at most $B_f\|v-v'\|_\infty$. Their denominators
are at least $\kappa$, and a numerator divided by its own denominator
belongs to $[0,1]$. Subtracting the two quotients gives a state
Lipschitz bound $2B_f/\kappa$. For fixed posterior, the difference of
report factors at two commands is at most $J\|g-g'\|_\infty$
pointwise, so the command bound is $2J/\kappa$. Consequently
\eqref{eq:compiler-lipschitz} holds uniformly on the calibration set.
Physical query predictions are fixed linear combinations of raw moments,
with uniformly bounded coefficients. For effective evaluation we fix
a sufficiently small positive rational $\delta$ in the canonical
probe basis $1/2$, $1/2+\delta t^{a_i}$. Its formal query coefficients
are then rational. Alternatively certified approximations to a chosen
fixed query coefficient table are included among the numerical inputs.
This verifies every metric
hypothesis of Theorem~\ref{thm:compact-compiler}.

\begin{lemma}[Certified finite evaluation from prior moments]
\label{lem:finite-moments}
For a fixed horizon, dyadic approximations to the coefficients of the
cell matrix and to
\begin{equation}\label{eq:finite-prior-table}
 m_\gamma(a)=\int_0^1t^{\gamma\cdot a}\,d\mu(t),\qquad
 \gamma\in\mathbb N_0^r,\quad |\gamma|=N,
\end{equation}
suffice to evaluate every rational-command history state and every fixed
finite product-query prediction with a prescribed error $\tau>0$.
Required input accuracy is $c_N\min\{\tau,1\}$ for a positive
constant determined by supplied coefficient and evidence bounds. The
number of input moment entries is independent of the history grid and
of $M$. No equality decision on the exponents is required.
\end{lemma}
\begin{proof}
Expand an $n$-report product in formal homogeneous variables as
$P(X)=\sum_{|\beta|=n}p_\beta X^\beta$. Its coefficient sum satisfies
$\sum_\beta|p_\beta|\le B_f^n$. A coordinate of the remaining raw
state is
\begin{equation}\label{eq:finite-evaluation}
 v_\alpha=
 \frac{\sum_{|\beta|=n}p_\beta m_{\alpha+\beta}}
      {\sum_{|\beta|=n}p_\beta m_{\beta+(N-n)e_0}},
 \qquad |\alpha|=N-n.
\end{equation}
The denominator is the true prefix evidence and is at least $\kappa^n$.
The formula remains exact when distinct formal labels have equal
exponents. There is no cancellation-based inference of rank.

A perturbation of size $\delta$ in the input coefficients changes each
factor coefficient vector by at most a fixed multiple of $\delta$.
The telescoping identity for products bounds the coefficient-sum error
in $P$ by $n(B_f+1)^{n-1}C\delta$, for $\delta$ sufficiently small.
Each prior moment lies in $[0,1]$, and there are finitely many formal
labels. Thus numerator and denominator errors are at most $D_N\delta$,
where $D_N$ is obtained from the finite bounds just displayed.
Take $D_N\delta\le\kappa^N/2$. The approximate denominator is then
positive. If a true numerator is $A$, its denominator is $Z$, and
$0\le A\le Z$, direct subtraction gives
\[
 \left|\frac{\widetilde A}{\widetilde Z}-\frac AZ\right|
 \le\frac{|\widetilde A-A|+|\widetilde Z-Z|}{\widetilde Z}
 \le4D_N\kappa^{-N}\delta.
\]
Choose $\delta$ to make this at most $\tau/2$ and round the result
to a dyadic within $\tau/2$. Fixed product-query predictions are
linear combinations of these raw entries, or the same product ratios;
increase the accuracy constant by their fixed coefficient bounds.
This gives the claimed finite procedure.
\end{proof}

\begin{theorem}[Finite-precision realization of the intrinsic law]
\label{thm:digital-intrinsic}
Fix the experiment, the full-support prior, the horizon and compact
calibration set of Section~\ref{sec:collision-geometry}. Suppose the
command cube is rationally specified, and positive rational lower bounds
for the report likelihoods and positive exponents, together with upper
coefficient bounds, are supplied. For each known calibration $a$ and
integer $M\ge1$, finite approximations to the data in
Lemma~\ref{lem:finite-moments} produce a deterministic finite-table
transducer with at most $M$ persistent labels per stage and
\begin{equation}\label{eq:digital-intrinsic}
 \sup_{n,\mathfrak h}\mathcal R(\widehat Q_n,Q_n(\mathfrak h))
 \le C\{\Xi_N(M,a)+h^2+\tau^2\}.
\end{equation}
Here $h$ is the command mesh and $\tau$ the certified error of the
offline state and output evaluations. Constants are uniform in $a,M$.
For meshes and evaluation errors at most $c/(M+1)$, this transducer
has regret at most $C'\Xi_N(M,a)$, while every such transducer with the memoryless input interface has
maximum unconditional checkpoint regret at least $c'\Xi_N(M,a)$
under the same exploration law as in the intrinsic theorem.

All numerical input and output precisions can be chosen
$\log_2(M+1)+O(1)$ binary places. The transition and output program has
$O(M^{J+1}\log(M+1))$ bits for the fixed experiment. Its synthesis
uses polynomially many fixed-dimensional rational operations in $M$,
with exponent allowed to depend on $J,N$. If the coefficients and
moments have computable approximation procedures, the synthesis is a
single uniform terminating algorithm using those procedures offline.
No computability assumption is needed for the preceding existential
intrinsic law for arbitrary full-support priors.
\end{theorem}
\begin{proof}
Theorem~\ref{thm:intrinsic-checkpoint}, its reachable-centre construction,
and the fixed physical/raw metric comparison imply
\[
 e_n(M)\le C_n\Psi_{n,N-n}(M,a)^{1/2}.
\]
The raw update calculation above and Lemma~\ref{lem:finite-moments}
verify all assumptions of Theorem~\ref{thm:compact-compiler}. Squaring
\eqref{eq:compiler-bound}, using
$(u+v+w)^2\le3(u^2+v^2+w^2)$, and averaging over the finite query
menu prove \eqref{eq:digital-intrinsic}. For $N\ge2$, the $\ell=1$
term at every nontrivial checkpoint gives
\[
 \Xi_N(M,a)\ge M^{-2}.
\]
Hence command and evaluation errors of order $(M+1)^{-1}$ are absorbed
without determining any individual pivot or collision scale.

A dyadic grid with this mesh has $O(M^J)$ points. Coefficient and prior
moment errors of size $c_N/(M+1)$ suffice by
Lemma~\ref{lem:finite-moments}; the required binary precision is the
stated logarithmic quantity. Proposition~\ref{prop:compiler-size} gives
the storage and arithmetic counts. Dyadic input coefficients may be
converted to rational product ratios of $O_N(\log(M+1))$ bits, and all
stored approximations have that precision. Integer comparisons are
finite, including exact ties.

Finally the finite-table filter, even when its command input is a
quantized approximation, is an admissible filter with at most $M$
labels in the original experiment when input coding is memoryless.
The exact-input filter can first apply that fixed coding rule. It has less, not more, information
than one reading the exact command. The lower bound in
Theorem~\ref{thm:intrinsic-streaming} therefore applies unchanged.
This is a lower bound for one filter over the common exploration law,
not a combination of independently redesigned checkpoint filters.
\end{proof}

\begin{corollary}[Calibration accuracy independent of collision gaps]
\label{cor:calibration-precision}
Suppose $a_i\ge a_*>0$ for $i>0$, the calibration entries are
computable reals, and the prior has a procedure that approximates
$\int t^b\,d\mu$ for positive rational $b$ to requested accuracy.
The preceding compiler needs only $\log_2(M+1)+O(1)$ binary places
of the calibration and of the resulting moments. The bound does not
involve the smallest nonzero additive exponent gap and is valid at
exact collisions.
\end{corollary}
\begin{proof}
For $b,b'\ge a_*/2$, the mean value theorem and
$\sup_{0<t\le1}t^{a_*/2}|\log t|=2/(e a_*)$ give
\[
 \left|\int(t^b-t^{b'})\,d\mu\right|
                   \le\frac{2}{e a_*}|b-b'|.
\]
Every nonconstant exponent in \eqref{eq:finite-prior-table} is at least
$a_*$. Approximate its constituent calibration entries so that its
error is $O((M+1)^{-1})$ and below $a_*/2$. Evaluate the corresponding
positive rational exponent to the same accuracy. The constant entry
is exactly one. There are finitely many labels, so the constants can
be chosen uniformly. This argument neither compares two sums nor asks
whether their difference vanishes.
\end{proof}

\begin{corollary}[Digital realization of high-contact phases]
\label{cor:digital-contact}
For the two-parameter family in Corollary~\ref{cor:two-parameter}, the
same three-term regret law is attained by finite numerical transition
and output tables under the finite-data assumptions above. Along
$u=\theta$, $v=\theta+\theta^k$, the required numerical precision is
$\log_2(M+1)+O(1)$, with no extra term proportional to
$\log(1/|v-u|)$. In particular the compiler remains valid before that
smallest separation can be numerically resolved.
\end{corollary}
\begin{proof}
Apply Theorem~\ref{thm:digital-intrinsic} and
Corollary~\ref{cor:calibration-precision}. The three-term expression and
its contact orders are those of Corollary~\ref{cor:two-parameter}.
The compiler never identifies the three collision lines or computes
$k$; its accuracy is set by $M$, and its performance follows from the
complete attainable cover. The physical experiment and its full-history
loss baseline are unchanged.
\end{proof}

The construction is not a new interpolation theorem or a new
farthest-first approximation algorithm. It connects a collision-uniform
attainability theorem to a finite numerical state machine with a counted
program. The distinction matters: a purely spectral identity supplies
neither the attainable cover nor the positive-likelihood update bound.
The horizon is fixed throughout. The program can be much larger than
the persistent state, and the running time of a general moment oracle
has not been bounded. For a noncomputable prior, finite advice exists
at each precision but a uniform algorithm for producing that advice is
not asserted. These distinctions leave the arbitrary-prior mathematical
classification intact.


\subsubsection{Complete numerical realization statement}
The following collects the intrinsic and effective conclusions. Its construction and resource bounds are proved in the surrounding appendices; no stronger numerical converse is implicit.
\begin{theorem}[Sharp causal memory and certified numerical realization]
\label{thm:effective-main}
Under the preceding fixed-experiment hypotheses, with $r\ge2$, the
checkpoint and causal regret laws have orders $\Psi_{n,m}$ and $\Xi_N$,
respectively, uniformly in $a\in K$ and integers $M\ge1$.
Suppose the command cube is rationally specified and certified finite
approximations to the cell coefficients and prior moments through formal
degree $N$ are available at requested tolerances, together with coefficient and
positivity bounds.
An algorithm constructs integer transition tables and dyadic output
tables with at most $M$ persistent labels at each stage and maximum
worst-history regret at most $C\Xi_N(M,a)$. For every such filter with
a fixed memoryless command-coding rule, the maximum unconditional
checkpoint regret under the common exploration law is at least
$c\Xi_N(M,a)$.

The algorithm requires no collision stratum, additive-gap lower bound,
or value of $\Xi_N$ as input. Its selected numerical precision satisfies
\[
 b(M,a)=\tfrac12\log_2(1/\Xi_N(M,a))+O(1),
\]
and a sufficient read-only program size is
$O(M\Xi_N(M,a)^{-J/2}\log(M+1))$ bits. In particular, setting
$d_0=(r-1)\lfloor N/2\rfloor$, one may uniformly use
\[
 b(M)\le d_0^{-1}\log_2(M+1)+O(1),\qquad
 P(M)=O(M^{1+J/d_0}\log(M+1)).
\]
These numerical bounds are sufficient; no precision or program-length
converse is asserted. The running filter retains only its index, not a
posterior vector, historical tape, or numerical oracle. For computable
numerical inputs the synthesis is uniform and terminating. All constants
may depend on the fixed prior and horizon, but not on $a$ or $M$.
\end{theorem}


\subsection{Certified precision and profile-adaptive compilation}
\label{sec:certified-resources}
The geometric profile has a second use besides specifying the error of
an $M$-state filter. It determines a sufficient numerical scale for
constructing that filter. We first obtain a scale uniform over the
calibration chamber. We then recover the actual covering scale from
finite approximate histories, without evaluating a determinant or
identifying a collision stratum. All precision and program-length
estimates in this section are upper resource bounds. The matching
converse remains the converse for persistent labels.

The separated-chain observation and its uniform resource consequence
were recorded in the preceding technical note \cite{A1v10note}.

\subsubsection{A separated family of future tests}
Set $k=\lfloor N/2\rfloor$ and
\[
 d_0=(r-1)k,\qquad
 \delta_K=\min_{a\in K}\min_{1\le i<r}(a_i-a_{i-1}),\qquad
 \bar\delta=\delta_K/H.
\]
We assume $r\ge2$, so $d_0\ge1$. Compactness in the strictly ordered
chamber gives $\delta_K>0$. This is a one-step separation, not a lower
bound on nonzero additive gaps.

\begin{proposition}[Separated-chain floor]\label{prop:separated-floor}
For every $a\in K$ and every integer $M\ge1$,
\begin{equation}\label{eq:separated-floor}
 \Xi_N(M,a)\ge \bar\delta^{\,d_0-1}M^{-2/d_0}.
\end{equation}
\end{proposition}
\begin{proof}
Put $m=N-k$. The future sumset $mA$ contains
\[
 C_m=\{jD:0\le j\le m\}
       \cup\{jD+a_i:0\le j<m,\ 1\le i\le r-2\}.
\]
Each displayed exponent is a sum of at most $m$ members of $A$;
padding by zeros produces a formal degree-$m$ label. In the interval
$[jD,(j+1)D]$ the chain is
\[
 jD<jD+a_1<\cdots<jD+a_{r-2}<(j+1)D.
\]
Thus $C_m$ has $m(r-1)+1$ distinct elements, with consecutive gaps at
least $\delta_K$. For $r=2$ there are no interior terms, and the same
count and separation hold. There are at least $d_0$ positive elements,
since $m\ge k$. In particular $q_m\ge d_0$ and
$p_{k,m}=d_0$.
Choose $d_0$ positive elements and one formal label for each. Every
pair of their normalized nodes is separated by at least $\bar\delta$.
Their Vandermonde product is at least
$\bar\delta^{d_0(d_0-1)/2}$. Keeping this one term in
$\Psi_{k,m}$ proves \eqref{eq:separated-floor}. Other determinant
terms, including all collision-sensitive ones, have not been removed.
\end{proof}

\begin{corollary}[Uniform sufficient numerical resources]
\label{cor:uniform-resources}
Under the finite-data hypotheses of Theorem~\ref{thm:digital-intrinsic},
the complete regret order $\Xi_N(M,a)$ is realized with at most $M$
persistent labels per stage using
\begin{equation}\label{eq:uniform-resources}
 b(M)=\left\lceil\frac{\log_2(M+1)}{d_0}\right\rceil+O(1),
 \qquad
 P(M)=O\bigl(M^{1+J/d_0}\log(M+1)\bigr)
\end{equation}
binary places and read-only program bits, respectively. With $T=N-1$,
a sufficient synthesis count is
\[
 O\bigl(M^{1+JT/d_0}+M^{2+J/d_0}\bigr)
\]
fixed-dimensional rational operations, excluding the running time of
numerical input procedures. The constants have the same fixed-experiment
and calibration-uniform scope as the intrinsic theorem.
\end{corollary}
\begin{proof}
Take $h,\tau$ bounded by a fixed multiple of $(M+1)^{-1/d_0}$.
Proposition~\ref{prop:separated-floor} gives
\[
 h^2+\tau^2\le C M^{-2/d_0}
              \le C\bar\delta^{1-d_0}\Xi_N(M,a).
\]
Theorem~\ref{thm:digital-intrinsic} therefore gives the full upper
law, not merely its separated-chain term. Its converse for memoryless
input coding is unchanged. Lemma~\ref{lem:finite-moments} transfers
this tolerance to the finitely many coefficient and moment inputs with
only a fixed additive change of precision. The grid has
$O(M^{J/d_0})$ points. Counting transition entries, output entries,
candidates, and nearest-centre comparisons as in
Proposition~\ref{prop:compiler-size} proves the remaining statements.
The variable part of $b$ is the least integer $j$ for which
$2^{j d_0}\ge M+1$, and can be found by integer comparisons. Neither
$\delta_K$ nor an additive-gap inverse is required to choose it.
\end{proof}

\subsubsection{An observable certificate for the covering scale}
Return to the general compact systems of Section~\ref{sec:effective}.
Write $e(M)=\max_{1\le n\le T}e_n(M)$ and choose a supplied rational
upper bound $A\ge\max_n A_n$. Use coordinate product grids with a fixed indexing convention, so their
mesh exponent is finite metadata rather than a list of real commands.
The constants $L_n,G_n,K_n$ may likewise
be replaced by supplied rational upper bounds. For a command $h$-net,
let $\widetilde S_n$ be the indexed list of approximate grid-history
states, each within $\tau$ of the corresponding true state. Let $I_n$
be the indices chosen by farthest-first selection with budget $M$, and
put
\begin{equation}\label{eq:observable-radius}
 r_n=\max_{z\in\widetilde S_n}\min_{i\in I_n}
                   \|z-\widetilde s_{n,i}\|_\infty,
 \qquad r=\max_{1\le n\le T}r_n.
\end{equation}
These numbers are rational and obtained from the finite data. The true
states of selected histories are denoted by $c_{n,i}$, only in the
analysis and in offline verification.

\begin{lemma}[Two-sided finite certificate]\label{lem:radius-certificate}
For each noninitial stage, and without any positivity assumption on
$e_n(M)$,
\begin{equation}\label{eq:radius-certificate}
 \max\{0,r_n/2-\tau\}\le e_n(M)
                   \le r_n+A_nh+2\tau.
\end{equation}
The selected true states cover $S_n$ with radius at most the right-hand
side. In particular, if $\tau=h$ and $c=A+2$, then
\begin{equation}\label{eq:max-radius-certificate}
 \max\{0,r/2-h\}\le e(M)\le r+ch.
\end{equation}
\end{lemma}
\begin{proof}
If the approximate list has more than $M$ entries, the selected $M$
points together with a farthest point have mutual distances at least
$r_n$. A cover of the true set by $M$ balls of radius $e_n(M)+\epsilon$
covers this list by balls of radius $e_n(M)+\epsilon+\tau$.
Two of the $M+1$ points belong to one ball, hence
$r_n\le2e_n(M)+2\epsilon+2\tau$. Let $\epsilon$ decrease to zero.
If there are at most $M$ entries, $r_n=0$ and the lower inequality is
immediate. Duplicate vectors cause no change in either argument.
For the upper inequality, approximate an arbitrary history by a grid
history. Its state changes by at most $A_nh$. Moving to the approximate
sample, then to a selected approximate centre, and then to its true
state costs at most $\tau+r_n+\tau$. This proves the covering claim
and hence the upper inequality. Taking maxima proves
\eqref{eq:max-radius-certificate}.
\end{proof}

\begin{theorem}[A posteriori precision selection]
\label{thm:adaptive-compiler}
Assume the compact-system and finite-evaluation hypotheses of
Theorem~\ref{thm:compact-compiler}, and suppose $e(M)>0$.
For $b=0,1,\ldots$, take $h_b=\tau_b=2^{-b}$, construct the lists
above, and stop at the first $b$ satisfying
\begin{equation}\label{eq:adaptive-stop}
 r_b\ge4c h_b,\qquad c=A+2.
\end{equation}
Every comparison is rational. The procedure terminates without knowing
$e(M)$. At termination it produces an $M$-label integer/dyadic program
with all-history query error at most $C_T e(M)$. Moreover
\begin{equation}\label{eq:adaptive-scale}
 h_b\asymp\min\{1,e(M)\},\qquad
 b=\log_2(2+e(M)^{-1})+O(1),
\end{equation}
where the constants depend only on the fixed system bounds. The final
finite certificate satisfies
\begin{equation}\label{eq:adaptive-bracket}
 \frac{3r_b}{8}\le e(M)\le\frac{5r_b}{4}.
\end{equation}
A sufficient read-only program size is
\begin{equation}\label{eq:adaptive-program}
 O\!\left(M(1+e(M)^{-1})^J\log(M+1)
                +M\log(2+e(M)^{-1})\right).
\end{equation}
\end{theorem}
\begin{proof}
By \eqref{eq:max-radius-certificate}, $r_b\ge e(M)-ch_b$.
Consequently the test succeeds whenever $h_b\le e(M)/(5c)$;
such a $b$ exists. At the first successful stage,
$h_b\le r_b/(4c)\le r_b/8$, since $c\ge2$.
Substituting in the lower and upper bounds of
\eqref{eq:max-radius-certificate} proves \eqref{eq:adaptive-bracket}.
It also gives $h_b\le e(M)/(2c-1)$, using
$r_b\le2e(M)+2h_b$.

If $b>0$, failure at $b-1$ and $h_{b-1}=2h_b$ imply
\[
 e(M)\le r_{b-1}+ch_{b-1}<5ch_{b-1}=10ch_b.
\]
Thus $e(M)/(10c)<h_b\le e(M)/(2c-1)$ in this case. If $b=0$,
then $h_b=1$ and the stopping condition and lower certificate give
$e(M)\ge2c-1$. Since $e(M)$ is also bounded by a constant depending
on $B$, the asserted comparison with $\min\{1,e(M)\}$ follows in
all cases. Taking logarithms proves the precision statement.

Use the final selected histories in the transition construction of
Theorem~\ref{thm:compact-compiler}. Equation~\eqref{eq:compiler-bound}
and $h_b=\tau_b\le C e(M)$ give the all-history error bound. The
number of state indices is not increased during this construction:
refinement stages are offline calculations, not persistent states.
The final command grid has $O((1+e(M)^{-1})^J)$ points. The transition
and output counts in Proposition~\ref{prop:compiler-size}, with
$b+O(1)$ places for outputs, give \eqref{eq:adaptive-program}.
No comparison of exact state coordinates, singular values, or exponents
occurs in the stopping rule.
\end{proof}

For $J\ge1$, the total count over all dyadic refinements is bounded,
up to fixed constants, by the last refinement count. Indeed the candidate
and transition bounds are geometric sums in $2^{bJ}$. A sufficient count
is
\[
 O\!\left(M(1+e(M)^{-1})^{JT}
                  +M^2(1+e(M)^{-1})^J\right).
\]
This counts rational operations, not oracle running time or unit-cost
operations on unrestricted real numbers. The number of refinements is
$O(\log(2+e(M)^{-1}))$. For fixed $M$ the procedure does not decide
whether an arbitrary compact set has zero covering radius.

\begin{corollary}[An absolute-tolerance version]
\label{cor:absolute-tolerance}
For any rational $\zeta\in(0,1]$, stop instead when either
\eqref{eq:adaptive-stop} holds or $h_b\le\zeta$. No assumption that
$e(M)>0$ is needed. The resulting all-history query error is at most
$C_T(e(M)+\zeta)$, and $h_b\asymp\min\{1,e(M)+\zeta\}$.
\end{corollary}
\begin{proof}
Termination follows from $h_b\downarrow0$. If the radius condition
holds, the previous theorem's stopping-stage estimates apply; if it
holds at $b>0$, failure of both tests at $b-1$ also gives
$h_b>\zeta/2$. If the radius condition fails at termination, the
certificate gives $e(M)<5ch_b$ and the tolerance condition gives
$h_b\le\zeta$; for $b>0$ one again has $h_b>\zeta/2$.
The initial-stage case is bounded by the fixed state diameter. In all
cases $h_b\asymp\min\{1,e(M)+\zeta\}$. Apply
\eqref{eq:compiler-bound} with $\tau_b=h_b$.
\end{proof}

\subsubsection{Verification of an individual finite program}
A small covering radius does not check that a stored decoder is correct.
The following residual certificate checks transitions and outputs
separately. It is useful also when a compilation uses sharper, local
numerical bounds than the universal covering estimate.

\begin{proposition}[Transition and decoder residual certificate]
\label{prop:program-certificate}
Consider any integer transition table $j(n,i,u,x)$ and any stored
probability vectors $q_{n,i}$. Suppose certified evaluations give
$\widetilde c_{n,i}$, $\widetilde y_{n,i,u,x}$ and
$\widetilde Q_{n,i}$ within $\tau$ of, respectively,
$c_{n,i}$, $T_{n,x}(c_{n,i},\bar u)$ and $Q_n(c_{n,i})$.
Define the rational residual bounds
\begin{align*}
 d_n&=\max_{i,u,x}
  \|\widetilde y_{n,i,u,x}-\widetilde c_{n+1,j(n,i,u,x)}\|_\infty
       +2\tau,\\
 o_n&=\max_i\|q_{n,i}-\widetilde Q_{n,i}\|_\infty+\tau.
\end{align*}
Put $D_0=0$ and $D_{n+1}=L_nD_n+G_nh+d_n$. For every actual
history and every admitted current input name,
\begin{equation}\label{eq:program-certificate}
 \|\widehat Q_n-Q_n(s_n)\|_\infty\le K_nD_n+o_n.
\end{equation}
The formula applies to a perturbed or corrupted table as well as to a
correctly compiled one; the measured residuals, and therefore the
certificate, then change.
\end{proposition}
\begin{proof}
Two state evaluation errors give
$\|T_{n,x}(c_{n,i},\bar u)-c_{n+1,j}\|_\infty\le d_n$.
One output evaluation error gives
$\|q_{n,i}-Q_n(c_{n,i})\|_\infty\le o_n$.
The transition Lipschitz bound and $\|u-\bar u\|_\infty\le h$
therefore imply inductively
$\|s_n-c_{n,i_n}\|_\infty\le D_n$. The query Lipschitz bound and
the output inequality prove \eqref{eq:program-certificate}. All
maxima are over finite tables; the justification for their numerical
error allowances is the supplied evaluation certificate, not an
assumption that approximate vectors are reachable states.
\end{proof}

\subsubsection{Finite input errors without consistency tests}
We record explicit constants for the finite moment interface. They are
not chosen to optimize numerical conditioning. Their purpose is to
certify the data actually used by the offline algorithm.

\begin{lemma}[A finite coefficient--moment certificate]
\label{lem:explicit-moment-certificate}
Suppose each true report factor has coefficient $\ell^1$ norm at most
$B\ge1$, and the coefficient matrix of $J$ cells and $r$ monomials
is supplied with entrywise error at most $\epsilon$. Every prior moment
through formal degree $N$ is supplied with error at most $\epsilon$.
Retain all formal labels, even when their exponents coincide. Put
$C_f=Jr$ and assume $C_f\epsilon\le1$. For a prefix of length $n$,
set
\[
 D_n^{\rm in}=nC_f(B+1)^{n-1}+(B+1)^n,
\]
where the first term is zero when $n=0$. The computed prefix evidence
$\widetilde Z$ and every computed raw-moment numerator
$\widetilde A$ have absolute errors at most
$\epsilon D_n^{\rm in}$. If
$\epsilon D_n^{\rm in}\le\kappa^n/2$, then
\begin{equation}\label{eq:quotient-certificate}
 \widetilde Z\ge\kappa^n/2>0,\qquad
 \left|\frac{\widetilde A}{\widetilde Z}-\frac A Z\right|
       \le\frac{2\epsilon D_n^{\rm in}}{\widetilde Z}
       \le4\epsilon D_n^{\rm in}\kappa^{-n}.
\end{equation}
For a true probability query with exact formal coefficient $\ell^1$
norm at most $B_Q$, replace the numerator error by
$B_Q\epsilon D_n^{\rm in}$ and the first numerator $2$ on the
right by $B_Q+1$.
\end{lemma}
\begin{proof}
For an accepted cell or a rejection mixture, command entries have
absolute value at most one. The error in a report factor has
$\ell^1$ norm at most $Jr\epsilon$. Thus every approximate factor
has norm at most $B+1$. Telescoping the product of $n$ factors in the
coefficient algebra, whose $\ell^1$ norm is submultiplicative, bounds
the product error by
$nC_f\epsilon(B+1)^{n-1}$. Integrating the product error against true
moments costs no further factor, since all true moments lie in $[0,1]$.
Integrating the approximate product against the moment errors costs at
most $\epsilon(B+1)^n$. Shifting its formal indices for a raw moment
changes neither coefficient norm. This proves both error estimates.
Since $Z\ge\kappa^n$, the asserted positivity follows. Finally,
with $q=A/Z\in[0,1]$,
\[
 \frac{\widetilde A}{\widetilde Z}-q
 =\frac{(\widetilde A-A)-q(\widetilde Z-Z)}{\widetilde Z}.
\]
The two numerator errors prove \eqref{eq:quotient-certificate}.
Multiplication by an exact query polynomial costs its coefficient norm,
and the same identity proves the query version. Clipping an approximate
probability to $[0,1]$ does not increase its error.
\end{proof}

The approximation table is not required to satisfy positivity identities
or equality between coincident formal moments. Those properties hold for
the true experiment; the error certificate is what is required of its
numerical data. Formal products of degree less than $N$ can be padded
by the zero exponent to degree $N$. Conversely, keeping all degrees
through $N$ gives the same finite interface. For a prescribed evaluation
tolerance $\tau$, choosing $\epsilon\le c_N\tau$ and allowing a fixed
number of additional dyadic places proves the required certificate.
Thus none of the refinement steps depends on deciding exact equality
of exponent sums.

\subsubsection{Recovery of the entire collision profile}
\begin{lemma}[Covering and prediction scales]\label{lem:cover-profile}
For the raw attainable state sets of the monomial experiment,
\begin{equation}\label{eq:cover-profile}
 e_n(M)^2\asymp\Psi_{n,N-n}(M,a),\qquad
 e(M)^2\asymp\Xi_N(M,a),
\end{equation}
uniformly over $a\in K$ and $M\ge1$.
\end{lemma}
\begin{proof}
The global reachable-centre construction in
Theorem~\ref{thm:intrinsic-checkpoint} gives the upper inequality.
For the reverse inequality take any $M$ raw centres with covering radius
$e_n(M)+\epsilon$. The fixed formal query matrix is a bounded linear
map from raw coordinates to physical prediction coordinates. Send the
index of a nearest raw centre and apply this matrix at the decoder;
clip each resulting probability to $[0,1]$. The worst-history mean
squared error is at most $C(e_n(M)+\epsilon)^2$. This is an admissible
one-checkpoint code, irrespective of whether its raw centres are
reachable. The worst-history lower bound in
Theorem~\ref{thm:intrinsic-checkpoint} gives
$c\Psi_{n,N-n}\le C(e_n(M)+\epsilon)^2$.
Let $\epsilon$ decrease to zero. There are finitely many checkpoints,
so taking maxima proves the second assertion with uniform constants.
This argument uses neither the new stopping rule nor its causal upper
bound, and hence does not assume the conclusion to be proved.
\end{proof}

\begin{theorem}[Profile-adaptive finite realization]
\label{thm:profile-adaptive}
Under the finite-data hypotheses of Theorem~\ref{thm:digital-intrinsic},
with numerical data available at successively requested tolerances,
one algorithm produces an $M$-label finite numerical filter without
receiving $\Xi_N$, a collision stratum, or any additive-gap lower bound.
Its maximum worst-history regret is at most $C\Xi_N(M,a)$, and the
original common-exploration lower bound is $c\Xi_N(M,a)$ for every
such filter with a memoryless input interface. The algorithm selects a
precision satisfying
\begin{equation}\label{eq:profile-precision}
 b(M,a)=\tfrac12\log_2\bigl(1/\Xi_N(M,a)\bigr)+O(1),
\end{equation}
and has sufficient read-only program size
\begin{equation}\label{eq:profile-program}
 P(M,a)=O\!\left(M\,\Xi_N(M,a)^{-J/2}\log(M+1)\right).
\end{equation}
Equation~\eqref{eq:profile-precision} describes the precision selected
by this algorithm, up to a uniform additive constant. It is not a
lower bound on the precision of every possible realization.
\end{theorem}
\begin{proof}
By Lemma~\ref{lem:cover-profile} and
Proposition~\ref{prop:separated-floor}, $e(M)>0$. All constants and
finite evaluation hypotheses for the raw system were verified in
Section~\ref{sec:effective}; Lemma~\ref{lem:explicit-moment-certificate}
provides explicit error allowances at each refinement. Apply
Theorem~\ref{thm:adaptive-compiler} and square its all-history query
bound. Equation~\eqref{eq:cover-profile} gives the desired upper law.
The state budget remains exactly the same upper budget $M$, so
memoryless input coding is simulated by an exact-command filter and
the intrinsic streaming converse applies unchanged.

Raw states lie in a unit cube, hence $e(M)\le1/2$. Also
$M^{-2}\le\Xi_N(M,a)\le1$. Taking logarithms of
\eqref{eq:cover-profile} and using \eqref{eq:adaptive-scale} proves
\eqref{eq:profile-precision}. Equations~\eqref{eq:adaptive-program}
and \eqref{eq:separated-floor} absorb the output term into
\eqref{eq:profile-program}, and give
$b(M,a)\le d_0^{-1}\log_2(M+1)+O(1)$.
The supplied dyadic inputs and rounded outputs have $b+O(1)$-place
accuracy. The exhibited exact rational evaluator uses $O(b+1)$-bit
intermediates, with a multiplicative constant depending on the fixed
horizon and input format; it need not have only an additive overhead.
Lemma~\ref{lem:rational-workspace} makes this distinction explicit. If moments are obtained from approximate calibration
entries, the derivative estimate in
Corollary~\ref{cor:calibration-precision} transfers the same tolerance.
The synthesis is uniform when those input procedures are computable;
otherwise it is a finite-advice construction with the same error bound.
\end{proof}

\begin{corollary}[Program-size phases at intersecting collision lines]
\label{cor:resource-phases}
In the four-cell, five-trial family of Corollary~\ref{cor:two-parameter},
write $\rho=\max\{|u|,|v|\}$ and
$\tau_c=\min\{|u|,|v-u|,|v-2u|\}$. Put
\[
 F=\max\{M^{-1/3},\ \rho^{1/2}M^{-1/4},\
                         (\rho^2\tau_c)^{2/9}M^{-2/9}\}.
\]
The adaptive algorithm realizes regret of order $F$ using
$b=\frac12\log_2(1/F)+O(1)$ places. A sufficient program bound is
\begin{equation}\label{eq:contact-program}
 O\!\left(\min\{M^{5/3},\ \rho^{-1}M^{3/2},\
                 (\rho^2\tau_c)^{-4/9}M^{13/9}\}\log(M+1)\right),
\end{equation}
where a reciprocal term with zero denominator is omitted. In particular
$b\le\frac16\log_2(M+1)+O(1)$ throughout the parameter square.
\end{corollary}
\begin{proof}
Here $r=J=4$, $N=5$, and $d_0=6$. The intrinsic two-parameter theorem
gives $\Xi_N\asymp F$. Substitute this comparison in
\eqref{eq:profile-precision} and \eqref{eq:profile-program}. The identity
$M F^{-2}=\min\{M^{5/3},\rho^{-1}M^{3/2},
(\rho^2\tau_c)^{-4/9}M^{13/9}\}$ proves
\eqref{eq:contact-program}, with the stated zero convention.
The separated-chain term gives the uniform precision bound.
\end{proof}

Along $u=\theta$, $v=\theta+\theta^k$, $k\ge2$, the three sufficient
program scales are respectively $M^{5/3}$,
$\theta^{-1}M^{3/2}$, and
$\theta^{-(4k+8)/9}M^{13/9}$, times $\log(M+1)$.
Their crossovers are the same $M\asymp\theta^{-6}$ and
$M\asymp\theta^{-(8k-2)}$ as for the error profile. This is a
consequence for numerical realizations, not a new classification of
collision arrangements or a converse for program length. In particular,
the algorithm need not determine $k$ or resolve $|v-u|$ before selecting
its numerical scale.


\subsection{Verified construction and stability of finite advice}
\label{sec:construction-stability}
The intrinsic law concerns an experiment, whereas a numerical certificate
concerns specified data or a specified program. We now relate these
statements without identifying them. The resulting construction is stable
under a common, imperfect numerical description of several experiments.
In particular, an exact additive relation need not be resolved when its
effect lies below the accuracy of the supplied moments.

\subsubsection{A construction certificate independent of the program residual}
Use the compact finite-horizon system of Section~\ref{sec:effective}.
The numerical value attached to each grid history is fixed during one
compilation. Let $\widetilde s_n(w)$ approximate its true state to accuracy
$\tau$ in the maximum norm, including numerical rounding. Denote the
selected histories by $w_{n,1},\ldots,w_{n,k_n}$, $k_n\le M$, and put
\[
 \widetilde c_{n,i}=\widetilde s_n(w_{n,i}),\qquad
 r_n=\max_{w\in\mathcal H_n^h}\min_{i\le k_n}
             \|\widetilde s_n(w)-\widetilde c_{n,i}\|_\infty.
\]
Here $\mathcal H_n^h$ is the complete finite list of grid histories.
A history extending a selected history by a grid command and a report
belongs to the next such list. We order the lists once and use the first
index whenever a minimum or maximum is attained more than once.

\begin{definition}[Finite construction contract]\label{def:construction-contract}
A table satisfies the construction contract if its representatives are
the farthest-first selection from the fixed numerical lists, its state
counts agree with those representatives, every transition satisfies
\begin{equation}\label{eq:nearest-contract}
 j(n,i,u,x)=\min\operatorname*{argmin}_{k\le k_{n+1}}
 \|\widetilde s_{n+1}(w_{n,i}(u,x))-
                               \widetilde c_{n+1,k}\|_\infty,
\end{equation}
and every output is the prescribed clipped dyadic rounding of the
numerical query evaluation at its representative history. All numerical
lists and their radii $r_n$ are fixed before a table is checked.
The histories and the contract are offline objects, not persistent state.
\end{definition}

\begin{theorem}[Construction-to-covering implication]
\label{thm:construction-contract}
Suppose the fixed state and query evaluations have errors at most $\tau$,
the command grid has radius $h$, and a table satisfies
Definition~\ref{def:construction-contract}. If $E_n$ bounds the distance
from a true state to the true representative selected by the running
filter, then one may take
\begin{equation}\label{eq:contract-recurrence}
 E_0=0,\qquad E_{n+1}=L_nE_n+G_nh+r_{n+1}+2\tau.
\end{equation}
The query error at stage $n$ is at most $K_nE_n+\tau$.
Moreover,
\begin{equation}\label{eq:contract-order}
 r_n\le 2e_n(M)+2\tau,\qquad
 \max_{i,u,x}\|T_n(c_{n,i},\bar u,x)-c_{n+1,j(n,i,u,x)}\|_\infty
       \le r_{n+1}+2\tau.
\end{equation}
The right side of the second inequality is fixed by the numerical
history lists, not increased in response to a supplied transition table.
Consequently, at fixed horizon the table has worst-history query error
$O(e(M)+h+\tau)$.
\end{theorem}
\begin{proof}
The true initial state is the representative of the empty history, so
$E_0=0$ even when its numerical name is imperfect. For a stored
representative and a grid command, the numerical target in
\eqref{eq:nearest-contract} belongs to $\mathcal H_{n+1}^h$.
Its distance to the selected numerical centre is therefore at most
$r_{n+1}$. Passing from the two numerical values to their two true
values costs at most $2\tau$. This proves the second inequality in
\eqref{eq:contract-order}. The first is the farthest-first separation
argument of Lemma~\ref{lem:radius-certificate}: if there is an
unselected point, the selected points together with a farthest point
are $r_n$-separated; transferring them to an arbitrary $M$-ball cover
of the true set gives $r_n\le2e_n(M)+2\tau$. If every point is
selected, $r_n=0$ and the inequality is immediate.

For an actual command, first compare the true current state with its
true representative, at cost $L_nE_n$, and then replace the command
by its grid code, at cost $G_nh$. Apply the just proved bound to the
appended representative history. This gives
\eqref{eq:contract-recurrence}. Compare a true query with its value
at that representative, at cost $K_nE_n$, and then use the certified
stored output, at cost $\tau$. Iterating the recurrence finitely many
times and using the first inequality in \eqref{eq:contract-order}
proves the last assertion. At no step is an approximate vector passed
to a transition map whose domain is the feasible state set.
\end{proof}

The residual bound in Proposition~\ref{prop:program-certificate} remains
valid for an arbitrary table. It is a different assertion: its right
side is measured from that table. The following comparison makes the
logical distinction quantitative. The example and its role as a
transition diagnostic were given in \cite{RefereeV11Note}; we include
the argument because the distinction is part of the numerical theorem.

\begin{proposition}[Residual validity and covering-order realization]
\label{prop:residual-comparison}
There is a fixed compact one-step Lipschitz system for which a sequence
of $M$-label tables has exact outputs at all representatives and valid
residual bounds, yet its worst-history squared prediction error divided
by the squared optimal $M$-centre radius is $4M^2$. Thus these two
properties, even together with knowledge of the optimal covering radius,
do not imply the construction-to-covering conclusion. For the same
system a table satisfying the nearest-centre construction has the
covering-order bound of Theorem~\ref{thm:construction-contract}.
\end{proposition}
\begin{proof}
Take $U=[1/4,3/4]$, $s_0=1/2$, $T(s,u)=(s+u)/2$, and $Q(s)=s$.
The reachable interval $[3/8,5/8]$ has length $1/4$. Since $M$ balls
of radius $e$ have total length at most $2Me$, its covering radius
is at least $1/(8M)$; equally spaced midpoints attain that value.
Select representatives by farthest-first on an endpoint-containing
command grid whose first history gives $3/8$, and store their exact
query values. Replace every transition target by this first index.
The output for command $3/4$ is now $3/8$, instead of $5/8$.
The maximum error is exactly $1/4$, so the asserted squared ratio is
$(1/4)^2/(1/(8M))^2=4M^2$. Its measured transition residual is also
$1/4$, hence its residual bound is true. None of these computations
changes the optimal radius of the experiment. For the correct table,
\eqref{eq:nearest-contract} and Theorem~\ref{thm:construction-contract}
give error $O(e(M)+h)$; choose $h=O(e(M))$. For the dyadic grids with
$8M$ segments used in the finite diagnostic, $M$ a power of two,
farthest-first gives maximum grid error at most $2e(M)$: after the
endpoints, each new point bisects a longest interval between consecutive
centres. Induction over the dyadic levels bounds its largest half-gap
by $1/(4M)$; for $M=1$ the endpoint error is $1/4$.
\end{proof}

This proposition is not a limitation on finite approximation methods.
It states exactly what a residual estimate does not assert about an
arbitrary program. The additional input that identifies the sharp
scale in this paper is the attained geometry of
Theorem~\ref{thm:intrinsic-checkpoint}, not the error recurrence itself.
We next give a consequence of that geometry which also concerns
imperfectly known experiments, rather than only a count of table entries.

\subsubsection{Stability of attainable sets and their profiles}
All state coordinates in this subsection are the formal raw coordinates;
coincident labels are not removed. Two experiments have the same
command and report alphabets and are compared on the same histories.
Their moment tables contain every formal degree through $N$.

\begin{lemma}[Paired-history stability]\label{lem:paired-history}
Consider two positive monomial experiments with uniformly bounded cell
coefficient norms and report factors bounded below by $\kappa>0$.
Suppose the maximum entrywise difference of their cell coefficients and
formal moment tables is at most $\delta\le1$. There is a constant $C_*$,
depending only on the fixed finite format and positivity bounds, such
that, at every history of length $n<N$,
\begin{equation}\label{eq:paired-stability}
 \|s_n(w)-s'_n(w)\|_\infty\le C_*\delta.
\end{equation}
For a common probe menu with fixed formal coefficients the query
probabilities satisfy the same assertion, with another fixed constant.
In particular,
\begin{equation}\label{eq:hausdorff-profile}
 d_H(S_n,S'_n)\le C_*\delta,\qquad
 |e_n(M)-e'_n(M)|\le C_*\delta\quad(M\ge1).
\end{equation}
No lower bound on a nonzero difference of exponent sums occurs here.
\end{lemma}
\begin{proof}
Let $B$ bound the coefficient norm of a report factor and put $C_f=Jr$.
For a fixed command, the difference of the two factor coefficient
vectors has norm at most $C_f\delta$. Telescoping a product of $n$
factors and then changing its moments bounds each integrated numerator
and its evidence by
\[
 D_n\delta,\qquad
 D_n=nC_f(B+1)^{n-1}+(B+1)^n.
\]
The same calculation applies after multiplication by any formal raw
monomial of the remaining degree. If $A/Z$ and $A'/Z'$ are the two
raw moments, then $Z,Z'\ge\kappa^n$ and $0\le A'/Z'\le1$.
Consequently
\[
 \left|\frac A Z-\frac{A'}{Z'}\right|
 =\frac{|A-A'+(A'/Z')(Z'-Z)|}{Z}
 \le 2D_n\kappa^{-n}\delta.
\]
Take the maximum over the finitely many stages. A fixed formal query
is a bounded linear combination of raw moments, proving its assertion.
These estimates compare two actual experiments; the formal moment
entries need not satisfy the same coincidence identities in the two.

Every reachable point is obtained from some history. Use that identical
history in the other experiment to obtain \eqref{eq:paired-stability}
in both directions. This proves the Hausdorff bound. An arbitrary
$M$-centre cover of one set, with radius within $\epsilon$ of its
infimum, covers the other with additional radius at most $C_*\delta$.
Interchange the sets and let $\epsilon$ decrease to zero to finish.
\end{proof}

Uniformity over priors must be specified, not inferred from full support
alone. Fix a full-support Borel probability $\mu_0$ on $[0,1]$ and
constants $0<c_-\le1\le c_+<\infty$. Write
\[
 \mathcal D=\{\mu\hbox{ a probability}:c_-\mu_0\le\mu\le c_+\mu_0\}.
\]
These are inequalities of measures. They do not require $\mu_0$, or
any member of $\mathcal D$, to have a Lebesgue density.

\begin{lemma}[Uniformity on a dominated prior family]
\label{lem:prior-envelope}
The constants in the intrinsic checkpoint and causal laws, and in
Lemma~\ref{lem:cover-profile}, can be chosen uniformly for
$(a,\mu)\in K\times\mathcal D$. The finite-horizon transition and
finite-input estimates have the same uniformity.
\end{lemma}
\begin{proof}
The family $\mathcal D$ is weakly compact. Indeed, probabilities on
$[0,1]$ are weakly compact, and each measure inequality is closed:
pass to the limit against a nonnegative continuous function and then
use regularity of finite Borel measures. Every member has full support.
The ordinary and complete confluent moment entries used in
Lemma~\ref{lem:confluent-positive} depend continuously on
$(a,\mu)$ on this compact family. For a logarithmic entry, its positive
exponent is bounded away from zero; hence $t^c(\log t)^j$, extended
by zero at the origin, varies continuously in the uniform norm on
the relevant compact parameter set. The zero-exponent constant block
has multiplicity one.

There are finitely many formal orderings and confluent block types.
For each, the mixed-moment pairing is strictly positive by the
confluent positivity lemma for every $\mu\in\mathcal D$.
Equivalently, in its determinant integration formula a product of
$k$ copies of $\mu$ dominates $c_-^k\mu_0^{\otimes k}$; the ordered
separated intervals used to prove strict positivity therefore retain
positive mass. Compactness and continuity give the uniform lower
singular-value bounds in the normalized attainable flag argument.
The derivatives used in its local inverse and volume minorization
are uniformly bounded, since all factor coefficients are bounded and
the evidence is at least $\kappa^n$. The acquisition-prefix probability
has this same lower bound independently of the prior.

For the global cover, changing the moments only changes real
coefficients in the bounded-format formula. Neither its format nor
its attainable dimension changes. The bounding rectangle, evidence
bound, and finite query maps are uniformly bounded. Thus both the
local lower estimate and the global upper estimate used in
Theorem~\ref{thm:intrinsic-checkpoint} are uniform. The causal argument
uses only these estimates and finitely many uniform transition bounds.
Finally, the proof of Lemma~\ref{lem:cover-profile} and the explicit
moment error estimate use those same constants. This accounts for
all prior-dependent inputs of the asserted estimates.
\end{proof}

\begin{corollary}[Stability of the collision profile]
\label{cor:profile-stability}
Fix the detector and let $(a,\mu),(a',\mu')\in K\times\mathcal D$.
If their formal moment tables differ by at most $\delta$, then,
uniformly for all $M\ge1$,
\begin{equation}\label{eq:xi-stability}
 \Xi_N(M,a)\le C\bigl(\Xi_N(M,a')+\delta^2\bigr),\qquad
 \Xi_N(M,a')\le C\bigl(\Xi_N(M,a)+\delta^2\bigr).
\end{equation}
For a fixed prior, $\delta$ may be replaced by $C\|a-a'\|_\infty$.
In particular, equality of additive sums may change without producing
an instability above the scale of the moment error.
\end{corollary}
\begin{proof}
Lemmas~\ref{lem:paired-history} and \ref{lem:prior-envelope} give
$e(M)\le e'(M)+C_*\delta$ and the reverse inequality, uniformly.
Square these inequalities and use Lemma~\ref{lem:cover-profile}.
For the calibration assertion, the constant formal label is independent
of $a$. For every other label of degree at most $N$, its exponent is
uniformly positive on $K$. The derivatives of $t^{\alpha\cdot a}$
with respect to positive-exponent coordinates are
$\alpha_jt^{\alpha\cdot a}\log t$, uniformly bounded on $[0,1]$.
The line segment between two points of $K$ retains a positive lower
bound on the one-step exponents. Integrating the derivative along that
segment proves the stated moment estimate for any probability prior.
\end{proof}

\subsubsection{One program from a common imperfect numerical name}
Suppose a collection of compact systems shares its labelled history
interface, its state-coordinate dimensions, and constants $L_n,G_n,K_n$.
A \emph{common numerical name of floor $\rho$}, $0<\rho\le1$, supplies
at scale $h=2^{-b}$ one finite command grid and one numerical evaluation
on its feasible histories. For every compatible system, the grid is an
$h$-net and its state and query errors, including rounding, are at most
\[
 \tau_b=\max\{h,\rho\}.
\]
The same data are used for all compatible systems. This quantifier is
essential: a different exact prior hidden in a program is not a common
numerical name. Write $e_{\mathcal E}(M)$ for the largest covering radius
of a compatible system $\mathcal E$, and let $A=\max_n A_n$ be the
history-net constant. Assume the state sets lie in fixed bounded cubes.

\begin{theorem}[Stable common-advice realization]\label{thm:common-advice}
From such a common name one can construct one table satisfying
Definition~\ref{def:construction-contract}, with at most $M$ persistent
labels at each stage, whose worst-history query error in every compatible
system is at most
\begin{equation}\label{eq:common-advice-error}
 C\bigl(e_{\mathcal E}(M)+\rho\bigr).
\end{equation}
The construction does not select a compatible system. More precisely,
put $c=A+4$, and stop at the first dyadic scale for which
\begin{equation}\label{eq:robust-stop}
 r_b\ge4c h_b\quad\hbox{or}\quad h_b\le\rho,
 \qquad r_b=\max_{1\le n<N}r_{n,b}.
\end{equation}
The stopping scale satisfies, simultaneously for every compatible system,
\[
 h_b\asymp\min\{1,e_{\mathcal E}(M)+\rho\}.
\]
All constants depend only on the fixed horizon, bounds, and dimensions.
\end{theorem}
\begin{proof}
The second condition ensures finite termination. Until and including the
first stopping scale, $\tau_b\le2h_b$: before crossing the floor it
is $h_b$, and the first dyadic scale below $\rho$ is greater than
$\rho/2$ (with equality at the initial endpoint treated separately).
For each compatible system, the finite-cover argument therefore gives
\begin{equation}\label{eq:robust-bracket}
 \max\{0,r_b/2-2h_b\}\le e_{\mathcal E}(M)
           \le r_b+(A+4)h_b.
\end{equation}
The numerical radius on both sides is common, although the true covering
radius may depend on the system.

Suppose $b>0$ is the first stopping index. At the previous scale both
conditions failed. By \eqref{eq:robust-bracket} and
$h_{b-1}=2h_b$,
\[
 e_{\mathcal E}(M)<10c h_b,\qquad \rho<2h_b.
\]
Thus $h_b>(e_{\mathcal E}(M)+\rho)/(10c+2)$.
If the radius condition holds at $b$, use
$r_b\le2e_{\mathcal E}(M)+4h_b$ to obtain
$h_b\le e_{\mathcal E}(M)/(2c-2)$.
If it does not, the floor condition gives $h_b\le\rho$.
In both cases $h_b\le e_{\mathcal E}(M)+\rho$.
The initial stopping case follows from the fixed diameter bound and
$\rho\le1$; if its radius condition holds, the same inequality bounds
$e_{\mathcal E}$ below by a fixed positive number, and otherwise
$\rho=1$. This proves the asserted scale comparison.

The farthest-first selection and the independently checked nearest
transitions use only the common numerical lists. In every compatible
system Theorem~\ref{thm:construction-contract} bounds the query error
by $C(e_{\mathcal E}+h_b+\tau_b)$. At the stopping scale the last
two terms are $O(e_{\mathcal E}+\rho)$, proving
\eqref{eq:common-advice-error}. Every state remains an index of an
actual representative history for that system. Neither its true
coordinates nor its model identity are consulted at runtime.
\end{proof}

\begin{corollary}[Finite moment uncertainty through collisions]
\label{cor:uncertain-monomial}
Fix a detector, $K$, $N$, and the prior envelope $\mathcal D$.
Suppose a finite numerical moment table of degree at most $N$ is within
$\delta$ of every compatible $(a,\mu)\in K\times\mathcal D$.
The cell coefficients may also be supplied to this numerical accuracy;
the underlying detector is fixed. For all sufficiently small $\delta>0$,
one causal $M$-label filter constructed from these same data has
maximum worst-history regret at most
\begin{equation}\label{eq:uncertain-profile}
 C\bigl(\Xi_N(M,a)+\delta^2\bigr)
\end{equation}
in every compatible experiment. No collision decision or value of $\Xi_N$
is supplied to the construction. Its selected scale is comparable to
$\min\{1,\sqrt{\Xi_N(M,a)}+\delta\}$ for each such experiment.
\end{corollary}
\begin{proof}
Lemma~\ref{lem:explicit-moment-certificate} bounds the error of evaluating
any feasible history from the supplied finite data by $C_0\delta$,
with $C_0$ uniform by Lemma~\ref{lem:prior-envelope}. Use
$\rho=2C_0\delta$, increasing $C_0$ if necessary to cover the fixed
query coefficients. Exact rational evaluation of the supplied entries,
followed by rounding at $b+2$ binary places, has error at most
$\rho/2+h_b/8\le\max\{h_b,\rho\}$. Every compatible experiment
therefore has the same valid numerical name. The rational command
cube admits the required grids. Apply Theorem~\ref{thm:common-advice},
square its bound, and use Lemma~\ref{lem:cover-profile}, uniformly in
the prior envelope. The physical probe map is the same fixed formal
map for the fixed detector. The running program has the prescribed
label budget throughout. The smallness restriction on $\delta$ merely
ensures the printed positive-evidence and $\rho\le1$ bounds.
\end{proof}

This is a robustness statement about one finite description, not an
assumption that the approximate moment table is itself a positive
moment sequence. It supplies a direct interpretation of the full
collision law when even exact identities among supplied entries are
unavailable. The additional $\delta^2$ term is not only an artifact of
propagating a norm bound.

\begin{proposition}[Sharp advice uncertainty in a positive experiment]
\label{prop:advice-floor}
There is a fixed positive detector with $r=2$ and $N=2$, and a dominated
full-support prior family, for which the minimax regret of one filter
required to serve all priors compatible with a common $\delta$-accurate
moment table has order
\[
 M^{-2}+\delta^2,
\]
for all $M\ge1$ and all sufficiently small $\delta>0$. The lower bound
holds under the same independent command exploration and a fixed finite
spanning probe menu. It applies to programs receiving that common advice;
it is not a precision converse for programs given different exact priors.
\end{proposition}
\begin{proof}
Take $a_0=0,a_1=1$ and the three detector cells
\[
 k_0(t)=\tfrac14,\qquad k_1(t)=\tfrac38+\tfrac18t,
 \qquad k_2(t)=\tfrac38-\tfrac18t.
\]
They are positive, sum to one, and span $\operatorname{span}\{1,t\}$.
Let $\eta=1/4$, and take as reference prior uniform measure on $[0,1]$,
with envelope constants $1/2$ and $3/2$. For
$\epsilon=6\delta\le1/2$ consider the two densities
$1\pm\epsilon(2t-1)$. Their moments are
\[
 m_k^\pm=\frac1{k+1}\pm\epsilon\frac{k}{(k+1)(k+2)}.
\]
Since $k/((k+1)(k+2))\le1/6$ for every nonnegative integer $k$,
the uniform-prior moments form one common $\delta$-accurate numerical
name, even if more moments than required are supplied.

Upon reporting accepted cell $0$, its likelihood $g_0/4$ is constant
in $t$, so the posterior remains the prior. The joint density of this
report and the observed command is identical in the two experiments;
its probability is at least $\eta/4$. The probe obtained with command
$(1/2,3/8,5/8)$ is the event of rejection, whose function is
$q(t)=1/2+t/32$. Its two posterior means differ by $\delta/16$.
For any possibly randomized prediction $Y$ using the common advice,
the distribution of $Y$ conditional on the accepted-cell-$0$ history
is the same in the two experiments. If the two means are $q_-,q_+$,
\[
 \frac{\E(Y-q_-)^2+\E(Y-q_+)^2}{2}
 =\E\left(Y-\frac{q_-+q_+}{2}\right)^2
       +\frac{(q_+-q_-)^2}{4}
 \ge\frac{\delta^2}{1024}.
\]
Integrate over those identical command histories and multiply by the
fixed positive probability of selecting this probe from the spanning
menu. The larger of the two unconditional regrets is at least
$c\delta^2$, independently of $M$.

For either fixed prior, the intrinsic checkpoint lower bound is
$c'M^{-2}$, uniformly in this envelope. Indeed, $n=m=1$ and the
attainable dimension is one, so $\Xi_2(M,a)=M^{-2}$.
A common-advice filter is an admissible $M$-label filter for that fixed
experiment; supplying less information does not invalidate its memory
lower bound. Combining the two lower bounds gives a constant times
$M^{-2}+\delta^2$. Corollary~\ref{cor:uncertain-monomial} gives the
matching upper bound simultaneously for all compatible priors in the
envelope, proving the assertion.
\end{proof}

\subsubsection{Exact rational workspace}
\begin{lemma}[Input accuracy and exact intermediate lengths]
\label{lem:rational-workspace}
For the finite moment evaluator at fixed horizon and fixed format,
$b+O(1)$-place dyadic inputs suffice for $O(2^{-b})$ evaluation error.
If evaluated with exact rational operations, every intermediate numerator
and denominator has $O(b+1)$ bits, with a constant depending on the fixed
arithmetic circuit. This statement does not give an additive $O(1)$
overhead in exact intermediate length. In fact, for
$x_b=(2^b-1)/2^b$, the exact fractional length of $x_b^N$ is $Nb$.
\end{lemma}
\begin{proof}
The input-accuracy assertion is the positive-evidence error bound of
Lemma~\ref{lem:explicit-moment-certificate}. For the exact length
assertion, bound the lengths of the numerator and denominator of a
reduced fraction by those of its unreduced representation. Multiplication
or division adds the operand lengths, while addition forms two products
and their sum, increasing the resulting length by at most one beyond
those product bounds. The circuit for one fixed-degree moment quotient
has a number of gates depending only on the fixed format. Induction over
these gates gives $C(b+1)$ bits for every intermediate. Its evidence is
nonzero by the supplied certificate. Distances and comparisons of a
fixed number of coordinates satisfy the same bound; selecting among
more histories increases the number of operations and the index length,
not the degree of a moment evaluation.

Finally, $2^b-1$ is odd. The reduced denominator of $x_b^N$ is thus
exactly $2^{Nb}$. This proves the stated exact length and explains why
input/output accuracy and exact arithmetic storage cannot be identified.
The rational-operation counts in the preceding sections exclude the
cost of obtaining numerical advice; they are not bit-operation counts.
\end{proof}


\subsection{Conformance to a finite synthesis request}
\label{sec:request-conformance}
The construction theorem assumes a specified horizon and a specified
total numerical error. A table cannot certify these assumptions by
choosing its own horizon or its own rounding scale. We record the finite
interface that connects the mathematical implication to the returned
program. This interface does not certify the accuracy of an arbitrary
external moment oracle.

\begin{definition}[Bound synthesis specification]
\label{def:bound-request}
A synthesis specification consists of the requested horizon $T$, command
and report alphabets with their fixed orders, state budget $M$, stage
coordinate dimensions $d_0,\ldots,d_T$, ordered query interfaces of
sizes $b_0,\ldots,b_T$, rounding precision $v$, raw evaluation error
bound $\sigma$, and total error allowance $\tau$, with
\begin{equation}\label{eq:bound-rounding-budget}
 \sigma+2^{-v-1}\le\tau.
\end{equation}
It also fixes the numerical state and query vectors on every grid history
through $T$. All these objects are fixed from the caller's task and
numerical input before synthesis is invoked. Their dimensions may be
specified directly by the application, or read from that external input
interface and checked on every history. They are never inferred from a
returned transition table. Repeated evaluations use this same frozen
name, with the same coordinate and query orders.
\end{definition}

An independently enumerated check first compares a returned program's
precision and all stage lengths with this specification. It then checks
the complete numerical lists, the farthest-first representatives, every
nearest-centre transition with the fixed tie rule, and every clipped
rounded output as in Definition~\ref{def:construction-contract}.
State and query dimensions are checked against the specification, not
merely against one another. The structural check precedes every use of
a program radius in a stopping certificate. Input fingerprints record
identity and provenance; they are not proofs of numerical accuracy.

\begin{proposition}[Request-bound construction implication]
\label{prop:request-conformance}
Suppose the numerical input fixed in Definition~\ref{def:bound-request}
has state and query error at most $\sigma$ on its specified histories,
and the grid is an $h$-net. If the preceding check accepts, then the
returned program implements all $T$ requested stages, uses at most $M$
labels per stage, and satisfies Theorem~\ref{thm:construction-contract}
with the caller's allowance $\tau$. Its independently recomputed radii
satisfy
\[
 \max\{0,r_n/2-\tau\}\le e_n(M)
      \le r_n+A_nh+2\tau.
\]
A returned object with a different precision or fewer stages is rejected
before these inequalities are issued. Changing a precision field without
changing its numerical tables does not avoid the entrywise checks.
\end{proposition}
\begin{proof}
The preliminary comparisons force the requested precision, stage lengths,
and interface dimensions. Rounding a real number to the nearest
$v$-place dyadic, with either prescribed rule at ties, changes it by at
most $2^{-v-1}$. The triangle inequality and
\eqref{eq:bound-rounding-budget} therefore give state and query errors
at most $\tau$; clipping query outputs to $[0,1]$ cannot increase
error to a true probability. The remaining finite checks are precisely
the hypotheses of Definition~\ref{def:construction-contract} on those
fixed numerical lists. Its theorem gives the recurrence and the label
budget. Lemma~\ref{lem:radius-certificate} gives the displayed bracket.
A mismatch of precision or of any stage length fails a preliminary
comparison. With correct metadata but different tables, acceptance still
requires every output and transition to equal the independently rebuilt
entry; if all entries agree, the program does implement the requested
finite construction. None of these implications uses a residual bound
increased after a faulty transition has been observed.
\end{proof}

For the ordinary and common-name adaptive constructions take $v=b+2$
and $\sigma=\tau_b/2$. Since $2^{-v-1}=h_b/8\le\tau_b/8$,
\eqref{eq:bound-rounding-budget} holds with the prescribed $\tau_b$.
The stopping bracket therefore uses an error allowance bound to the
request before synthesis. Freezing these finite inputs and checking them
is offline work. It adds no field to the running state, which remains
one integer index. The earlier residual certificate for an arbitrary
table remains valid, but does not assert conformance to a different
synthesis request.
